\subsection{\texorpdfstring{第~\ref{sec:differentiable-maps}~节}{可微映射}}\label{subsec:solutions-differentiable-maps}

\begin{description}[font=\normalfont,leftmargin=2em,labelwidth=1.5em]
\RGsolution{ex:smooth-map-is-continuous}{smooth-map-is-continuous}
给定 $N$ 中的开集 $W$, 要证明 $f^{-1}(W)\subset M$ 是开集. 由 Exercise~\ref{ex:exercise-differentiable-structures}(\ref{ex:open-sets-in-parameterizations}), 只需证明对 $M$ 的每个参数化 $(U,\varphi)$, 集合 $\widetilde U:=\varphi^{-1}(f^{-1}(W))\subset\mathbb{R}^m$ 是开集 (其中 $m=\dim M$). 考虑 $N$ 的一个图册 $\{(V_\beta,\psi_\beta)\}$, 有
\begin{align*}
  \widetilde U
  &=\varphi^{-1}(f^{-1}(W))
  =\varphi^{-1}\left(\bigcup_\beta f^{-1}(W\cap\psi_\beta(V_\beta))\right)\\
  &=\bigcup_\beta\varphi^{-1}\left(f^{-1}\left(\psi_\beta\left(\psi_\beta^{-1}(W)\right)\right)\right)\\
  &=\bigcup_\beta(\psi_\beta^{-1}\circ f\circ\varphi)^{-1}\left(\psi_\beta^{-1}(W)\right).
\end{align*}
由于 $W$ 是开集, $\psi_\beta^{-1}(W)$ 是 $\mathbb{R}^n$ 中的开集 ($n=\dim N$). 又因为映射 $\psi_\beta^{-1}\circ f\circ\varphi$ 是 $C^\infty$ 的, 因而连续, 所以集合 $(\psi_\beta^{-1}\circ f\circ\varphi)^{-1}(\psi_\beta^{-1}(W))$ 也是开集, 由此可知 $\widetilde U$ 是开集.

\RGsolution{ex:smooth-antipodal-map}{smooth-antipodal-map}
对 $S^n$ 使用图册 $\{(\mathbb{R}^n,\pi_N^{-1}),(\mathbb{R}^n,\pi_S^{-1})\}$ [参见 Exercise~\ref{ex:exercise-differentiable-structures}(\ref{ex:stereographic-sphere-atlas})]. 要检验 $f$ 可微, 必须证明 $\pi_N\circ f\circ\pi_N^{-1}$, $\pi_S\circ f\circ\pi_N^{-1}$, $\pi_N\circ f\circ\pi_S^{-1}$ 和 $\pi_S\circ f\circ\pi_S^{-1}$ 这四个映射都可微. 由于 $f(S)=N$, 第一个映射的定义域为
\begin{align*}
  \pi_N(f^{-1}(\pi_N^{-1}(\mathbb{R}^n)))
  &=\pi_N(f^{-1}(S^n\setminus\{N\}))\\
  &=\pi_N(S^n\setminus\{S\})=\mathbb{R}^n\setminus\{0\}
\end{align*}
(为使表述清楚, 这里稍微滥用了记号). 有
\begin{align*}
  \pi_N\circ f\circ\pi_N^{-1}(y)
  &=\pi_N\circ f\left(\frac{2y}{1+\lVert y\rVert^2},\frac{\lVert y\rVert^2-1}{\lVert y\rVert^2+1}\right)\\
  &=\pi_N\left(-\frac{2y}{1+\lVert y\rVert^2},\frac{1-\lVert y\rVert^2}{1+\lVert y\rVert^2}\right)
  =-\frac{y}{\lVert y\rVert^2},
\end{align*}
它在 $\mathbb{R}^n\setminus\{0\}$ 上可微. 同理可证, 其余三个映射都在各自的定义域上可微.

\RGsolution{ex:properties-of-differentiable-maps-part-six}{different-smooth-structures-on-the-line}
\begin{enumerate}[label=(\alph*),ref=\alph*,leftmargin=*]
\item 恒等映射不是微分同胚, 因为 $\widehat{\operatorname{id}}(x):=(\varphi_2^{-1}\circ\operatorname{id}\circ\varphi_1)(x)=x^{1/3}$ 在原点不可微.
\item 映射 $f$ 是微分同胚, 因为 $\widehat f:=\varphi_2^{-1}\circ f\circ\varphi_1=\operatorname{id}$ 和 $\widehat{f^{-1}}:=\varphi_1^{-1}\circ f^{-1}\circ\varphi_2=\operatorname{id}$ 都是 $C^\infty$ 的.
\end{enumerate}
\end{description}
